\chapter{Streszczenie}

Projekt zrealizowany w ramach przedmiotu „Projektowanie interfejsów użytkownika” to makieta aplikacji internetowej dla sieciowej pizzerii „Wito's”, która została pierwotnie przygotowana w programie Figma, a następnie zaimplementowana w czystym kodzie HTML i CSS. Wybrane widoki zostały również opracowane jako responsywne strony internetowe, co potwierdza ich adaptacyjność do ekranów desktopowych i mobilnych. Aplikacja stanowi dwojaki system, który od strony klienckiej umożliwia łatwe przeglądanie menu, układanie własnych dań w interaktywnym kreatorze pizzy oraz wygodną obsługę koszyka zakupowego. Dodatkowo zaimplementowano funkcjonalny portal użytkownika, pozwalający zarejestrowanym klientom na edycję profilu, przeglądanie historii swoich zamówień czy zgłaszanie ewentualnych zwrotów. Z drugiej strony, aplikacja zawiera widoki dedykowanego panelu administracyjnego, dzięki któremu pracownicy lokalu mogą w przejrzysty sposób zarządzać dostępnym, w restauracji asortymentem i samą kartą menu. Interfejs zoptymalizowano pod kątem nowoczesnych standardów estetycznych oraz intuicyjności nawigacji, co uwidacznia silne skupienie na jak najlepszych doświadczeniach końcowego użytkownika (UX). Całość projektu stanowi przez to kompletne, gotowe wizualnie środowisko e-commerce, które precyzyjnie symuluje przepływ pracy zarówno z perspektywy głodnego klienta, jak i sprzedawcy.

